\paragraph*{04.03.03.01. Die für das Projekt gültigen Rechtsvorschriften identifizieren und einhalten}
\kcilabel{para:04.03.03.01_Rechtsvorschriften}{04.03.03.01}
\label{para:04.03.03.01_Rechtsvorschriften}
%%%% Task
\par Als \gls{wsl} \gls{ricefw} identifizierte und bewertete ich relevante gesetzliche, regulatorische und konzerninterne Vorschriften (\gls{dsgvo}, \gls{dora}, \gls{kritis}, \gls{gob}/\gls{gobd}, Compliance, Sicherheitsrichtlinien, \gls{pcr}-Anforderungen SAP, luftfahrttechnische Regularien für \gls{adn}) und überführte sie in ein verbindliches Steuerungsmodell. Ziel: Auditierbares, reproduzierbares Compliance-Management über sieben Zulieferer.

%%% Action
\par Bestellte programmweites Rechts- und Standardregister, konsolidierte Regularien und ließ es durch \gls{jur}, \gls{com}, \gls{jdr} und \gls{itsec} validieren. Definierte Checkpoints mit regulatorischer Klassifizierung, Prüfungen (Archivierung, Logging), Clean-Core-Bewertungen und \gls{gobd}-Prüfpunkten.

%Result
\par Alle \gls{ricefw}-Objekte durchliefen Compliance-Check; 100\% Checks für Welle 1 vorhanden. Keine kritischen Verstöße in \gls{uat}, Cutover-Quality-Gate und \gls{wp}-Prüfung. \gls{pcr}-Qualität nicht erreicht wegen Wechsels zu \gls{gfa}; als \gls{td} dokumentiert, Refactoring abgestimmt.\endnote{Bewusst eingegangen \gls{td} wird durch nachträgliche Restrukturierung geheilt} Nachweis: Fähigkeit, regulatorische Anforderungen systematisch zu identifizieren und einzuhalten.

%Reflect
\par Zukünftig Einbezug von Rechts- und Compliancefunktionen in frühe Backlog-Refinements verlagern, um Anforderungen rechtssicher einzuordnen. 